%% generato da strumenti/estrai_livelli.py: non modificare a mano
\subsection{Luce, atmosfera e partenza}
\begin{itemize}
\item \nodo{Directional Light} (il sole): Mobility \emph{Movable}, Rotation pitch -32, yaw -35; Intensity 9 lux; \emph{Use Temperature} sì, Temperature 5200 K; \emph{Atmosphere Sun Light} sì; Source Angle 1.
\item \nodo{Sky Atmosphere} (valori di default).
\item \nodo{Sky Light}: Movable, \emph{Real Time Capture} sì, Intensity 1.
\item \nodo{Exponential Height Fog} a Z $-200$: Fog Density 0,01.
\item \nodo{Post Process Volume}: \emph{Infinite Extent (Unbound)} sì; spuntate e impostate Auto Exposure Bias 0,3, Bloom Intensity 0,5, Vignette Intensity 0,25.
\item \nodo{Player Start} in (-500; 0; 120), rotazione 0.
\end{itemize}

\subsection{Partenza e primo tratto}
\begin{oggetti}
\nota{acqua tutto intorno: chi cade perde una vita e torna all'ultimo posto sicuro (la zona sta SOPRA l'acqua, che ha collisione: si "affoga" prima di toccarla)}
1 & Cubo & Acqua & (4500; 0; -310) & (17000; 9000; 20) & MI\_\allowbreak{}AIV\_\allowbreak{}Acqua \\
2 & BP\_\allowbreak{}ZonaCaduta & Zona\_\allowbreak{}Caduta & (4500; 0; -200) & -- & ZoneExtent = (8500; 4500; 150); ShowVisual = no \\
\end{oggetti}

\subsection{Isola di partenza}
\begin{oggetti}
3 & Pedana & Isola\_\allowbreak{}Partenza & (0; 0; -150) & (1600; 1600; 300) & MI\_\allowbreak{}AIV\_\allowbreak{}Erba, si cammina a Z = 0 \\
4 & Text Render & Scritta & (-150; 0; 330) & yaw 180 & ``LIVELLO 1 / salta, raccogli, scopri'', World Size 55, colore (1; 1; 1) \\
5 & Text Render & Scritta & (-150; 0; 230) & yaw 180 & ``3 VITE: se cadi in acqua torni all'ultimo posto sicuro / (a zero vite si ricomincia dal livello 1)'', World Size 26, colore (1; 0,55; 0,65) \\
\nota{il faro: piu' ti allontani, piu' diventa grande (cosi' lo vedi sempre)}
6 & Cubo & Faro\_\allowbreak{}Palo & (-650; 650; 200) & (30; 30; 400) & MI\_\allowbreak{}AIV\_\allowbreak{}LegnoScuro \\
7 & BP\_\allowbreak{}CresceConDistanza & Faro\_\allowbreak{}CresceConLaDistanza & (-650; 650; 450) & -- & ShapeMesh = Sphere; Material = MI\_\allowbreak{}AIV\_\allowbreak{}Abilita; BaseSize = (100; 100; 100); MinScale = 1; MaxScale = 7; NearDistance = 400; FarDistance = 9000; ScaleHeight = sì; Label = ``FARO: piu' sei lontano, piu' cresce'' \\
8 & BP\_\allowbreak{}Moneta & Moneta\_\allowbreak{}Isola\_\allowbreak{}1 & (-300; -300; 90) & -- & Value = 1; Material = MI\_\allowbreak{}AIV\_\allowbreak{}Moneta; GlowIntensity = 4; GlowColor = (1; 0,75; 0,3) \\
9 & BP\_\allowbreak{}Moneta & Moneta\_\allowbreak{}Isola\_\allowbreak{}2 & (-50; -300; 90) & -- & Value = 1; Material = MI\_\allowbreak{}AIV\_\allowbreak{}Moneta; GlowIntensity = 4; GlowColor = (1; 0,75; 0,3) \\
10 & BP\_\allowbreak{}Moneta & Moneta\_\allowbreak{}Isola\_\allowbreak{}3 & (200; -300; 90) & -- & Value = 1; Material = MI\_\allowbreak{}AIV\_\allowbreak{}Moneta; GlowIntensity = 4; GlowColor = (1; 0,75; 0,3) \\
\nota{pozzanghera: il materiale fisico dell'acqua cambia il suono dei passi}
11 & Cubo & Pozzanghera & (-250; 350; 1) & (360; 260; 2) & MI\_\allowbreak{}AIV\_\allowbreak{}Acqua \\
12 & Decal & Pozza\_\allowbreak{}Bordo & (-250; 350; 10) & rot (-90; 0; 0), Decal Size (30; 140; 190) & MI\_\allowbreak{}AIV\_\allowbreak{}Pozza \\
\nota{mulino: le pale ruotano, la testa ti segue (con un limite all'angolo)}
13 & BP\_\allowbreak{}Mulino & Mulino & (-450; -650; 0) & yaw 60 & BladeSpeed = 50; LookAtPlayer = sì; YawLimit = 60; TowerHeight = 900; BladeLength = 420; BladesCollide = no; TowerMaterial = MI\_\allowbreak{}AIV\_\allowbreak{}IntonacoChiaro; BladeMaterial = MI\_\allowbreak{}AIV\_\allowbreak{}LegnoChiaro \\
\nota{casse fisiche: spawn con due cicli for e un trace verso il basso}
14 & BP\_\allowbreak{}CasseFisiche & Casse\_\allowbreak{}Fisiche & (250; 550; 300) & -- & Rows = 3; Columns = 3; CubeLayers = 2; CubeSize = 60; Spacing = 66; SimulatePhysics = sì; CubeMass = 6; Materials = [MI\_\allowbreak{}AIV\_\allowbreak{}LegnoChiaro, MI\_\allowbreak{}AIV\_\allowbreak{}LegnoScuro] \\
\nota{staccionata lungo una spline (un paletto ogni 200 cm)}
15 & BP\_\allowbreak{}Staccionata & Staccionata & (-780; -780; 0) & -- & Spacing = 200; Rails = 2; Material = MI\_\allowbreak{}AIV\_\allowbreak{}LegnoScuro; Collide = sì; punti della spline (locali): (0; 0; 0); (0; 1560; 0); (1560; 1560; 0) \\
16 & Static Mesh & Cespuglio\_\allowbreak{}0 & (-600; 620; 0) & scala (1,3; 1,3; 1,3) & SM\_\allowbreak{}Groundplant\_\allowbreak{}12; Collision: No Collision \\
17 & Static Mesh & Cespuglio\_\allowbreak{}1 & (500; -700; 0) & scala (1,3; 1,3; 1,3) yaw 70 & SM\_\allowbreak{}Groundplant\_\allowbreak{}12; Collision: No Collision \\
18 & Static Mesh & Cespuglio\_\allowbreak{}2 & (-700; 100; 0) & scala (1,3; 1,3; 1,3) yaw 140 & SM\_\allowbreak{}Groundplant\_\allowbreak{}12; Collision: No Collision \\
\nota{alberi sulle isolette intorno (foliage "a mano")}
19 & Pedana & Isoletta\_\allowbreak{}Nord & (1500; 2600; -150) & (2400; 1400; 300) & MI\_\allowbreak{}AIV\_\allowbreak{}Erba, si cammina a Z = 0 \\
20 & Pedana & Isoletta\_\allowbreak{}Sud & (4500; -2900; -150) & (3200; 1600; 300) & MI\_\allowbreak{}AIV\_\allowbreak{}Erba, si cammina a Z = 0 \\
21 & BP\_\allowbreak{}Spargi & Alberi\_\allowbreak{}Nord & (1500; 2600; 300) & -- & Mesh = SM\_\allowbreak{}Chestnut\_\allowbreak{}4; AreaExtent = (1100; 600; 600); Count = 5; Seed = 11; MinScale = 0,55; MaxScale = 0,8; MinDistance = 600; AlignToGround = sì; Collide = sì; CullDistance = 0 \\
22 & BP\_\allowbreak{}Spargi & Alberi\_\allowbreak{}Sud & (4500; -2900; 300) & -- & Mesh = SM\_\allowbreak{}Chestnut\_\allowbreak{}4; AreaExtent = (1500; 700; 600); Count = 7; Seed = 23; MinScale = 0,55; MaxScale = 0,8; MinDistance = 600; AlignToGround = sì; Collide = sì; CullDistance = 0 \\
\end{oggetti}

\subsection{Cancello: pedana a pressione $\rightarrow$ porta doppia + luce}
\begin{oggetti}
\nota{muro con un'apertura al centro}
23 & Cubo & Cancello\_\allowbreak{}Muro\_\allowbreak{}S & (780; -460; 200) & (60; 680; 400) & MI\_\allowbreak{}AIV\_\allowbreak{}Pietra \\
24 & Cubo & Cancello\_\allowbreak{}Muro\_\allowbreak{}N & (780; 460; 200) & (60; 680; 400) & MI\_\allowbreak{}AIV\_\allowbreak{}Pietra \\
25 & Cubo & Cancello\_\allowbreak{}Architrave & (780; 0; 350) & (60; 240; 100) & MI\_\allowbreak{}AIV\_\allowbreak{}Pietra \\
26 & BP\_\allowbreak{}Porta & Porta\_\allowbreak{}Doppia & (780; 0; 0) & -- & DoorWidth = 240; DoorHeight = 300; DoorThickness = 16; Sliding = no; OpenAngle = 100; TransitionTime = 1,2; DoorMaterial = MI\_\allowbreak{}AIV\_\allowbreak{}LegnoDipinto \\
27 & BP\_\allowbreak{}PiastraPressione & Pedana\_\allowbreak{}Pressione & (450; -300; 0) & -- & StayPressed = sì; OnMaterial = MI\_\allowbreak{}AIV\_\allowbreak{}Acceso; OffMaterial = MI\_\allowbreak{}AIV\_\allowbreak{}Spento; Targets = [Porta\_\allowbreak{}Doppia] \\
28 & BP\_\allowbreak{}LuceAttivabile & Lampada\_\allowbreak{}Cancello & (720; 0; 360) & -- & TransitionTime = 0,6; ListenTo = [Pedana\_\allowbreak{}Pressione] \\
29 & Text Render & Scritta & (450; -300; 140) & yaw 180 & ``sali sulla pedana'', World Size 35, colore (1; 1; 1) \\
\end{oggetti}

\subsection{Gradini}
\begin{oggetti}
30 & Pedana & P1 & (1100; 0; -20) & (400; 400; 40) & MI\_\allowbreak{}AIV\_\allowbreak{}Percorso, si cammina a Z = 0 \\
31 & Pedana & P2\_\allowbreak{}Gradino & (1550; 0; 40) & (300; 300; 40) & MI\_\allowbreak{}AIV\_\allowbreak{}PercorsoScuro, si cammina a Z = 60 \\
32 & Pedana & P3\_\allowbreak{}Gradino & (1900; 0; 100) & (300; 300; 40) & MI\_\allowbreak{}AIV\_\allowbreak{}Percorso, si cammina a Z = 120 \\
33 & Pedana & P4\_\allowbreak{}Gradino & (2300; 0; 160) & (300; 300; 40) & MI\_\allowbreak{}AIV\_\allowbreak{}PercorsoScuro, si cammina a Z = 180 \\
34 & BP\_\allowbreak{}Moneta & Moneta\_\allowbreak{}P2 & (1550; 0; 150) & -- & Value = 1; Material = MI\_\allowbreak{}AIV\_\allowbreak{}Moneta; GlowIntensity = 4; GlowColor = (1; 0,75; 0,3) \\
35 & BP\_\allowbreak{}Moneta & Moneta\_\allowbreak{}P3 & (1900; 0; 210) & -- & Value = 1; Material = MI\_\allowbreak{}AIV\_\allowbreak{}Moneta; GlowIntensity = 4; GlowColor = (1; 0,75; 0,3) \\
36 & BP\_\allowbreak{}Moneta & Moneta\_\allowbreak{}P4 & (2300; 0; 270) & -- & Value = 1; Material = MI\_\allowbreak{}AIV\_\allowbreak{}Moneta; GlowIntensity = 4; GlowColor = (1; 0,75; 0,3) \\
\end{oggetti}

\subsection{Piattaforma mobile (suono con pitch diverso all'andata e al ritorno)}
\begin{oggetti}
37 & BP\_\allowbreak{}PiattaformaMobile & Mobile\_\allowbreak{}1 & (2650; 0; 165) & -- & EndOffset = (600; 0; 0); TravelTime = 2,5; PauseTime = 0,8; StartDelay = 0; PlatformSize = (250; 250; 30); PlatformMaterial = MI\_\allowbreak{}AIV\_\allowbreak{}Salto; MoveSound = SfxPlatformMoving\_\allowbreak{}Cue \\
38 & Text Render & Scritta & (2450; 0; 420) & yaw 180 & ``aspetta la piattaforma'', World Size 35, colore (1; 1; 1) \\
\end{oggetti}

\subsection{Corridoio con la zona velocità e salto lungo}
\begin{oggetti}
39 & Pedana & P5\_\allowbreak{}Corridoio & (3900; 0; 160) & (900; 400; 40) & MI\_\allowbreak{}AIV\_\allowbreak{}Percorso, si cammina a Z = 180 \\
40 & BP\_\allowbreak{}ZonaModificatore & Zona\_\allowbreak{}Velocita & (3650; 0; 280) & -- & ZoneExtent = (180; 180; 100); WalkSpeedMultiplier = 1,8; DurationAfterExit = 4; Label = ``VELOCITA x1.8''; VisualMaterial = MI\_\allowbreak{}AIV\_\allowbreak{}Salto \\
41 & Text Render & Scritta & (4250; 0; 420) & yaw 180 & ``prendi la rincorsa!'', World Size 40, colore (1; 1; 1) \\
\end{oggetti}

\subsection{IL DOPPIO SALTO}
\begin{oggetti}
42 & Pedana & P6\_\allowbreak{}DoppioSalto & (5000; 0; 160) & (500; 500; 40) & MI\_\allowbreak{}AIV\_\allowbreak{}Salto, si cammina a Z = 180 \\
43 & BP\_\allowbreak{}PickupAbilita & Oggetto\_\allowbreak{}DoppioSalto & (5000; 0; 290) & -- & Ability = DoppioSalto; Material = MI\_\allowbreak{}AIV\_\allowbreak{}Abilita; GlowIntensity = 20; GlowColor = (0,3; 0,7; 1); PickupRadius = 80; Label = ``DOPPIO SALTO''; Message = ``Doppio salto! Premi salto di nuovo in aria: puoi anche cambiare direzione'' \\
44 & Spot Light & Luce\_\allowbreak{}DoppioSalto & (5000; 0; 900) & pitch -90 yaw 0 & 400 cd, raggio 1200, colore (0,5; 0,75; 1), coni 12/22° \\
\nota{parete di 3,3 m: con un salto solo (2,15 m) non ci si arriva}
45 & Pedana & P7\_\allowbreak{}Parete & (5550; 0; 345) & (500; 500; 330) & MI\_\allowbreak{}AIV\_\allowbreak{}PercorsoScuro, si cammina a Z = 510 \\
46 & Text Render & Scritta & (5250; 0; 560) & yaw 180 & ``troppo alta? salta due volte'', World Size 35, colore (1; 1; 1) \\
\nota{salti in cui, se sbagli, ti correggi in aria con il secondo salto}
47 & Pedana & P8 & (6250; 0; 490) & (300; 300; 40) & MI\_\allowbreak{}AIV\_\allowbreak{}Percorso, si cammina a Z = 510 \\
48 & Pedana & P9 & (6900; 200; 580) & (300; 300; 40) & MI\_\allowbreak{}AIV\_\allowbreak{}PercorsoScuro, si cammina a Z = 600 \\
49 & Pedana & P10 & (7550; -100; 490) & (300; 300; 40) & MI\_\allowbreak{}AIV\_\allowbreak{}Percorso, si cammina a Z = 510 \\
50 & BP\_\allowbreak{}Moneta & Moneta\_\allowbreak{}P8 & (6250; 0; 600) & -- & Value = 1; Material = MI\_\allowbreak{}AIV\_\allowbreak{}Moneta; GlowIntensity = 4; GlowColor = (1; 0,75; 0,3) \\
51 & BP\_\allowbreak{}Moneta & Moneta\_\allowbreak{}P9 & (6900; 200; 690) & -- & Value = 1; Material = MI\_\allowbreak{}AIV\_\allowbreak{}Moneta; GlowIntensity = 4; GlowColor = (1; 0,75; 0,3) \\
52 & BP\_\allowbreak{}Moneta & Moneta\_\allowbreak{}P10 & (7550; -100; 600) & -- & Value = 1; Material = MI\_\allowbreak{}AIV\_\allowbreak{}Moneta; GlowIntensity = 4; GlowColor = (1; 0,75; 0,3) \\
53 & Text Render & Scritta & (6250; 0; 780) & yaw 180 & ``se manchi il salto, risalta in aria / e correggi la direzione'', World Size 32, colore (1; 1; 1) \\
\end{oggetti}

\subsection{GENERATORE DI CASSE: la parete di 4,6 m}
\begin{oggetti}
54 & Pedana & P11\_\allowbreak{}Casse & (8300; 0; 490) & (700; 600; 40) & MI\_\allowbreak{}AIV\_\allowbreak{}Salto, si cammina a Z = 510 \\
55 & BP\_\allowbreak{}PickupAbilita & Oggetto\_\allowbreak{}Casse & (8100; 0; 620) & -- & Ability = GeneratoreCasse; Material = MI\_\allowbreak{}AIV\_\allowbreak{}Abilita; GlowIntensity = 20; GlowColor = (0,3; 0,7; 1); PickupRadius = 80; Label = ``GENERATORE DI CASSE''; Message = ``Generatore di casse! Premi F: fino a 3 casse, poi 5 secondi di ricarica'' \\
56 & Spot Light & Luce\_\allowbreak{}Oggetto\_\allowbreak{}Casse & (8100; 0; 1210) & pitch -90 yaw 0 & 300 cd, raggio 1000, colore (0,5; 0,75; 1), coni 12/22° \\
57 & Pedana & P12\_\allowbreak{}ParetePiuAlta & (8900; 0; 590) & (500; 600; 760) & MI\_\allowbreak{}AIV\_\allowbreak{}PercorsoScuro, si cammina a Z = 970 \\
58 & Text Render & Scritta & (8300; 0; 760) & yaw 180 & ``parete di 4,6 m: nemmeno il doppio salto basta / F davanti al muro: 2 casse una sull'altra, salici e salta due volte'', World Size 32, colore (1; 0,9; 0,5) \\
\end{oggetti}

\subsection{PONTE FORTUNATO (c'e' solo il 35\% delle volte) e la strada lunga}
\begin{oggetti}
59 & BP\_\allowbreak{}PresenzaCasuale & Ponte\_\allowbreak{}Fortunato & (9625; 0; 955) & -- & SpawnChance = 0,35; Size = (950; 250; 30); Material = MI\_\allowbreak{}AIV\_\allowbreak{}Moneta; Label = ``PONTE FORTUNATO (35\%)'' \\
60 & Pedana & Sasso\_\allowbreak{}1 & (9450; -450; 950) & (200; 200; 40) & MI\_\allowbreak{}AIV\_\allowbreak{}PercorsoScuro, si cammina a Z = 970 \\
61 & Pedana & Sasso\_\allowbreak{}2 & (9850; -450; 980) & (200; 200; 40) & MI\_\allowbreak{}AIV\_\allowbreak{}PercorsoScuro, si cammina a Z = 1000 \\
62 & Text Render & Scritta & (9150; 0; 1230) & yaw 180 & ``il ponte c'e' solo 35 volte su 100 / se non c'e': i sassi a sinistra'', World Size 32, colore (1; 0,9; 0,5) \\
\end{oggetti}

\subsection{Il totem che cambia modello ogni secondo}
\begin{oggetti}
63 & Pedana & P13 & (10400; 0; 950) & (600; 600; 40) & MI\_\allowbreak{}AIV\_\allowbreak{}Percorso, si cammina a Z = 970 \\
64 & BP\_\allowbreak{}CambiaModello & Totem\_\allowbreak{}CambiaModello & (10450; 220; 1040) & -- & Interval = 1; Size = (120; 120; 140); Collide = sì; Materials = [MI\_\allowbreak{}AIV\_\allowbreak{}Salto, MI\_\allowbreak{}AIV\_\allowbreak{}Abilita, MI\_\allowbreak{}AIV\_\allowbreak{}Moneta, MI\_\allowbreak{}AIV\_\allowbreak{}Pericolo]; Label = ``CAMBIA MODELLO OGNI SECONDO'' \\
65 & BP\_\allowbreak{}Moneta & Moneta\_\allowbreak{}P13 & (10350; -150; 1060) & -- & Value = 2; Material = MI\_\allowbreak{}AIV\_\allowbreak{}Moneta; GlowIntensity = 4; GlowColor = (1; 0,75; 0,3) \\
\end{oggetti}

\subsection{Traguardo}
\begin{oggetti}
66 & Pedana & P14\_\allowbreak{}Arrivo & (11150; 0; 950) & (600; 600; 40) & MI\_\allowbreak{}AIV\_\allowbreak{}Traguardo, si cammina a Z = 970 \\
67 & BP\_\allowbreak{}Traguardo & Traguardo & (11200; 0; 1070) & -- & ZoneExtent = (150; 150; 100); NextLevel = L\_\allowbreak{}GD13\_\allowbreak{}Livello2\_\allowbreak{}BP; Message = ``Livello 1 completato! Le monete sono al sicuro''; Delay = 1,2; ResetDeathCount = no; Label = ``TRAGUARDO''; VisualMaterial = MI\_\allowbreak{}AIV\_\allowbreak{}Traguardo \\
68 & Point Light & Luce\_\allowbreak{}Traguardo & (11200; 0; 1230) & -- & 30 cd, raggio 700, colore (0,4; 1; 0,5), niente ombre \\
\end{oggetti}

